% Figure: gradient-B drift. In a field that weakens from left to right the
% Larmor radius is larger where B is smaller, so successive gyro-loops do not
% close and the guiding centre marches sideways. Ions and electrons drift in
% opposite directions because the gyration sense depends on charge, so the
% drift carries a net current.
\begin{figure}[htbp]
\centering
\begin{tikzpicture}[scale=1.0]
  % field-strength indication: tightly packed dots on the left, sparse right
  \foreach \x in {0,0.45,...,6.0} {
    \foreach \y in {-1.4,-0.7,...,1.5} {
      \node[engblue, font=\tiny] at (\x,\y) {$\cdot$};
    }
  }
  \node[engblue, font=\scriptsize, anchor=south west] at (0,1.7)
    {strong $B$ (small $r_L$)};
  \node[engblue, font=\scriptsize, anchor=south east] at (6,1.7)
    {weak $B$ (large $r_L$)};
  \draw[engblue, ->] (0,-1.9) -- (6,-1.9)
    node[midway, below, font=\scriptsize] {decreasing $|\vect{B}|$, i.e.\ $\grad B$};
  % ion gyro-orbits of growing radius, guiding centre marching right (upper)
  \draw[engred, thick]
    (0.4,0.9) arc (180:-160:0.22)
    (1.1,0.95) arc (180:-150:0.30)
    (2.0,1.0) arc (180:-140:0.40)
    (3.1,1.05) arc (180:-130:0.52);
  \draw[engred, ->, thick] (4.0,0.9) -- (5.2,0.9);
  \node[engred, font=\scriptsize, anchor=south] at (4.6,0.95)
    {ion drift};
  % electron gyro-orbits, guiding centre marching left (lower)
  \draw[engorange, thick]
    (5.6,-0.9) arc (0:340:0.22)
    (4.9,-0.95) arc (0:330:0.30)
    (4.0,-1.0) arc (0:320:0.40);
  \draw[engorange, ->, thick] (2.9,-0.9) -- (1.7,-0.9);
  \node[engorange, font=\scriptsize, anchor=south] at (2.3,-0.85)
    {electron drift};
\end{tikzpicture}
\caption[Gradient-B drift]{The gradient-$B$ drift in a field that weakens
from left to right. A gyrating particle has a larger Larmor radius on the
weak-field part of its orbit than on the strong-field part, so the loop does
not close on itself and the guiding centre steps sideways, perpendicular to
both $\vect{B}$ and $\grad B$. The step size scales with $r_L$ and with the
fractional change in field over a Larmor radius. Ions (upper, red) and
electrons (lower, orange) gyrate in opposite senses, so they drift in
opposite directions, and the gradient drift therefore carries a net
current, in contrast to the charge-independent $E\times B$ drift. The same
charge separation, produced by gradient and curvature drifts in a toroidal
field, is what a tokamak must short out with its helical field twist to keep
the plasma confined.}
\label{fig:vol04ch13:gradb-drift}
\end{figure}
